% Appendix D - CLI reference.
% Source: the binary itself (--help) and the argparse definitions in
% src/univorch/interfaces/cli/app.py. When new commands or flags are
% added, propagate here.

\umaappendix{CLI reference}
\label{ap:cli-reference}

This appendix documents the \texttt{univorch} command-line client.
The CLI is a pure HTTP client of the daemon: every command translates
into one or more calls to the REST API and prints the result. Inside
the production Docker container the binary is on the \texttt{PATH};
on a development machine it is available with \texttt{uv run
univorch}.

\section{Invocation modes}

The CLI supports the same set of commands in two modes:

\begin{description}
  \item[Single-shot (bash mode)] -- one command per invocation, suitable
        for shell scripts.
        Example: \texttt{univorch deploy /lab/networks/student01}.
  \item[Interactive REPL] -- launched with no arguments. Carries a
        prompt that shows the current folder and (optionally) the
        remote daemon. Supports history, tab-completion of paths and
        in-shell \texttt{help <command>}.
\end{description}

\section{Selecting the daemon}

The CLI must know where the daemon listens. The selection follows
this precedence (first match wins):

\begin{enumerate}
  \item The \texttt{-{}-remote URL} argument on the command line.
  \item The \texttt{UNIVORCH\_REMOTE} environment variable.
  \item The built-in default \texttt{http://localhost:8080}.
\end{enumerate}

Inside the REPL, the \texttt{connect <URL>} command switches the
target daemon for the rest of the session and updates the prompt to
show the new target.

\section{Command reference}

\Cref{tab:cli-commands} lists the implemented commands at the close
of Sprint~4. The detailed description of each command follows.

\begin{table}[ht]
  \centering
  \small
  \begin{tabular}{l p{0.62\linewidth}}
    \toprule
    \textbf{Command} & \textbf{Purpose} \\
    \midrule
    \texttt{cd [path]}        & Change the current folder in the REPL.\\
    \texttt{pwd}              & Print the current folder.\\
    \texttt{ls [path]}        & List one level of a folder (alias of \texttt{list}).\\
    \texttt{list [path]}      & Same as \texttt{ls}.\\
    \texttt{tree [path] [-d]} & Print the full subtree under a folder.\\
    \texttt{inspect <path> [-l]} & Show the resolved (or local-only) definition.\\
    \texttt{status <path>}    & Show the descriptor state and the runtime state.\\
    \texttt{deploy <path>}    & Move a descriptor from \texttt{provisioned} to \texttt{deployed}.\\
    \texttt{undeploy <path>}  & Move a \texttt{deployed} descriptor back to \texttt{provisioned}.\\
    \texttt{start <path>}     & Power on a deployed VM.\\
    \texttt{stop <path>}      & Power off a deployed VM.\\
    \texttt{load <file> [dest]} & Load a YAML definition into a folder.\\
    \texttt{connect <url>}    & (REPL only) switch the live target daemon.\\
    \texttt{help [cmd]}       & Built-in help; one command without args lists all.\\
    \texttt{quit}             & (REPL only) leave the shell.\\
    \bottomrule
  \end{tabular}
  \caption{Commands provided by the \texttt{univorch} CLI.}
  \label{tab:cli-commands}
\end{table}

\subsection{Navigation: \texttt{cd}, \texttt{pwd}}

\texttt{cd [path]} changes the current folder of the REPL. The path
follows the conventions of a Unix shell: it may be absolute or
relative, \texttt{.} and \texttt{..} are honoured, and \texttt{..} at
the root stays at the root. With no argument, \texttt{cd} does
nothing; \texttt{pwd} prints the current folder.

\subsection{Listing: \texttt{ls}, \texttt{list}, \texttt{tree}}

\texttt{ls} (alias \texttt{list}) prints one level of children of a
folder, with a glyph that marks the state of each descriptor.
\texttt{tree} prints the entire subtree under a folder. The flag
\texttt{-d} (or \texttt{-{}-folders-only}) of \texttt{tree} hides
descriptors and shows only the folder structure, the equivalent of
the Linux \texttt{tree -d}.

The state glyphs used by both commands match the iconography of the
web GUI: an empty square for \texttt{provisioned}, a filled square
for \texttt{deployed}, a cross for \texttt{broken} and a warning
triangle for \texttt{unreachable}.

\subsection{Inspection: \texttt{inspect}, \texttt{status}}

\texttt{inspect <path>} shows the entity at a tree path. For
descriptors, the default mode merges the template, the cascade
inheritance and the lexical closure, so the displayed definition is
the \emph{effective} one. With \texttt{-l} (or \texttt{-{}-local}),
only the fields explicitly written at the node are shown.

\texttt{status <path>} reports both axes of the descriptor's state:
the orchestrator's state (\texttt{provisioned}, \texttt{deployed},
\texttt{broken}, \texttt{unreachable}) and, when applicable, the
runtime state read from the hypervisor (\texttt{running},
\texttt{stopped}, \texttt{paused}).

\subsection{Lifecycle: \texttt{deploy}, \texttt{undeploy}, \texttt{start}, \texttt{stop}}

The four lifecycle commands take a descriptor path and submit the
matching operation to the daemon. The daemon performs the action
through the connector that the descriptor's hypervisor declares.

\texttt{deploy} moves the descriptor from \texttt{provisioned} to
\texttt{deployed}. \texttt{undeploy} moves it back to
\texttt{provisioned} and removes the VM from the hypervisor.
\texttt{start} and \texttt{stop} change the runtime state of an
already deployed VM. All four are idempotent: ordering an action
that is already satisfied returns success with the message
\emph{``no change''} (DEC-035).

\subsection{Loading definitions: \texttt{load}}

\texttt{load <file> [dest]} reads a YAML definition file and
materialises it inside the destination folder. The file describes
the folders and descriptors to add; the destination is the parent.
Without a destination, the current folder is used. The YAML is
relative to the destination, which is what makes a definition file
portable across installations.

The load operation is fail-fast: if any item fails validation, the
whole load is rejected and the tree is not modified.

\subsection{REPL-only commands: \texttt{connect}, \texttt{help}, \texttt{quit}}

\texttt{connect <url>} switches the live session to a different
daemon. The prompt updates to show the new target between square
brackets (for example \texttt{univorch [http://aulario.uma.es:8080]
/lab/>}). The change is not persisted; the next REPL session goes
back to the precedence rule above.

\texttt{help} alone lists the commands; \texttt{help <cmd>} prints
the argparse-generated help for a specific command, with the same
syntax and examples shown above. \texttt{quit} (or \texttt{Ctrl+D}
on an empty line) leaves the shell.

\section{Examples}

A typical session inside the production container:

\begin{verbatim}
$ docker exec -it univorch univorch
univorch /> load /opt/univorch/examples/setup.yml /
6 items applied.
univorch /> cd /lab/networks
univorch /lab/networks/> ls
[ ] student01
[ ] student02
[ ] student03
univorch /lab/networks/> deploy student01
deploy /lab/networks/student01: ok (5.6s)
univorch /lab/networks/> status student01
descriptor state : deployed
runtime state    : stopped
vm id            : mock-7f3a91
univorch /lab/networks/> start student01
start /lab/networks/student01: ok
univorch /lab/networks/> quit
\end{verbatim}

The same operations from a shell script:

\begin{verbatim}
#!/bin/bash
export UNIVORCH_REMOTE=http://orquestador.dept.uma.es:8080
univorch load /opt/univorch/examples/setup.yml /
for s in student01 student02 student03; do
    univorch deploy /lab/networks/$s
done
\end{verbatim}

The CLI surface is intentionally small: the same vocabulary appears
in the YAML model, in the REST API, in the CLI and in the web GUI.
Adding new commands in the future means adding new operations to the
orchestrator's facade and exposing them; the CLI's argparse
boilerplate is the smaller part of the work.
